\documentclass[11pt, letterpaper]{article}
\input{../../../report.tex}
\usepackage{titlepageBU}
\title{A Proof of the Snake Lemma (For $R$-Modules)}
\date{February 7-8th, 2025}
\DeclareMathOperator{\im}{im}
\DeclareMathOperator{\coker}{coker}
\declaretheoremstyle[headfont=\bfseries, bodyfont=\normalfont, qed=\qedsymbol]{proof2}
\declaretheorem[numbered=no, style=proof2, name=Proof]{replacementproof2}
\renewenvironment{proof}[1][]{\begin{replacementproof2}}{\end{replacementproof2}}
\togglelogo
\begin{document}
\makereport
\section{Proof}
\begin{lemma}[The Snake Lemma]\label{snake}
	Let the diagram of $R$-modules
	\[\begin{tikzcd}[row sep = 2.5em, column sep = 2.5em]
		0\arrow[r]&L_1\arrow[r, "\alpha_1"]\arrow[d, "\lambda "]&M_1\arrow[r, "\beta _1"]\arrow[d, "\mu "]&N_1\arrow[r]\arrow[d, "\nu "]&0\\
		0\arrow[r]&L_0\arrow[r, "\alpha _0"]&M_0\arrow[r, "\beta _0"]&N_1\arrow[r]&0
	\end{tikzcd}\]
	commute, where rows are exact complexes. Then there exists an exact sequence
	\[\begin{tikzcd}[row sep = 2.5em, column sep = 2.5em]
		0\arrow[r]&\ker \lambda \arrow[r]&\ker \mu \arrow[r]&\ker \nu \arrow[r, "\delta "]&\coker \lambda \arrow[r]&\coker \mu \arrow[r]&\coker \nu \arrow[r]&0.
	\end{tikzcd}\]
\end{lemma}
\begin{proof}
	For convenience of notation, we will denote rows by $1_\bullet $ and $0_\bullet $, and columns by $L_\bullet $, $M_\bullet $, and $N_\bullet $.

Clearly, there exists the morphism $0 : 0 \to \ker \lambda $ with $\ker 0=0$ and $\im 0=\left\{ 0 \right\} $.

Since $1_\bullet $ is exact, $\ker \alpha_1=\im 0=\left\{ 0 \right\} $. Consider the restriction $\alpha_1|_{\ker \lambda }$. Note that for any $a\in \ker \lambda $,
\[\begin{tikzcd}[row sep = 2.5em, column sep = 2.5em]
	a\arrow[r, "\alpha_1"]\arrow[d, "\lambda "']&\alpha_1(a)\arrow[d, "\mu "]\\
	0\arrow[r]&0
\end{tikzcd}\]
because the diagram commutes. We must then have $\im \alpha_1|_{\ker \lambda }\subseteq \ker \mu $. Therefore, the complex
\[\begin{tikzcd}[row sep = 2.5em, column sep = 2.5em]
	0\arrow[r, "0"]&\ker \lambda \arrow[r, "\alpha_1|_{\ker \lambda }"]&\ker \mu 
\end{tikzcd}\]
is exact at $0,\ker \lambda $.

Now consider the restriction $\beta _1 |_{\ker \mu }: \ker \mu  \to N_1$. First, we will show $\im \beta_1\subseteq \ker \nu $. Let $a\in \ker \mu $, and note that
\[\begin{tikzcd}[row sep = 2.5em, column sep = 2.5em]
	a\arrow[r, "\beta_1"]\arrow[d, "\mu "']&\beta_1(a)\arrow[d, "\nu "]\\
	0\arrow[r]&0
\end{tikzcd}\]
because the diagram commutes. Therefore, $\beta_1|_{\ker \mu }: \ker \mu  \to \ker \nu $. Also note that because $1_\bullet $ is exact, we have $a\in \im \alpha_1 \iff a\in \ker \beta_1$. Restricting these sets to the subset $\ker \mu $ does not change this fact, so we have that the complex
\[\begin{tikzcd}[row sep = 2.5em, column sep = 2.5em]
	0\arrow[r, "0"]&\ker \lambda\arrow[r, "\alpha_1|_{\ker \lambda }"]&\ker \mu \arrow[r, "\beta_1|_{\ker \mu }"]&\ker \nu 
\end{tikzcd}\]
is exact at $\ker \mu $.

Now we must show the morphism $\delta : \ker \nu  \to \coker \lambda $ exists, which will be done by a diagram chase. Consider the diagram
\begin{equation}\label{eq:1}
	\begin{tikzcd}[row sep = 2.5em, column sep = 2.5em]
	&0\arrow[d]&0\arrow[d]&0\arrow[d]\\
	0\arrow[r]&\ker \lambda\arrow[r]\arrow[d]&\ker \mu \arrow[r]\arrow[d]&\ker \nu \arrow[d]\\
	0\arrow[r]&L_1\arrow[r, "\alpha_1"]\arrow[d, "\lambda "]&M_1\arrow[r, "\beta_1"]\arrow[d, "\mu "]&N_1\arrow[r]\arrow[d, "\nu "]&0\\
	0\arrow[r]&L_0\arrow[r, "\alpha _0"]\arrow[d]&M_0\arrow[r, "\beta _0"]\arrow[d]&N_0\arrow[r]\arrow[d]&0\\
		  &\coker \lambda \arrow[r]\arrow[d]&\coker \mu \arrow[r]\arrow[d]&\coker \nu \arrow[r]\arrow[d]&0\\
		  &0&0&0
\end{tikzcd},\end{equation}
where all unexplained arrows are induced completely naturally. Let $a\in \ker \nu $, and consider the chase
\begin{equation}\label{eq:2}
\begin{tikzcd}[row sep=2em, column sep=3em]
	&0\arrow[d, dotted, no head]&0\arrow[d, dotted, no head]&0\arrow[d, dotted, no head]\\
	0\arrow[r, dotted, no head]&\arrow[r, dotted, no head]\arrow[d, dotted, no head]&\arrow[d,  dotted, no head]\arrow[r, dotted, no head]&a\arrow[d]\\
	0\arrow[r, dotted, no head]&\arrow[r, dotted, no head]\arrow[d, dotted, no head]&c\arrow[r, "\beta _1"]\arrow[d, "\mu "]&b\arrow[r, dotted, no head]\arrow[d, dotted, no head, "\nu "]&0\\
	0\arrow[r, dotted, no head]&e\arrow[r, "\alpha _0"]\arrow[d]&d\arrow[d, dotted, no head]\arrow[r, dotted, no head, "\beta _0"']&*\arrow[r, dotted, no head]\arrow[d, dotted, no head]&0\\
				   &f\arrow[r, dotted, no head]\arrow[d, dotted, no head]&\arrow[r, dotted, no head]\arrow[d, dotted, no head]&\arrow[r, dotted, no head]\arrow[d, dotted, no head]&0\\
				   &0&0&0
\end{tikzcd}.
\end{equation}
We will show this diagram induces a homomorphism $\delta : \ker \nu  \to \coker \lambda $. The map $\ker \nu \hookrightarrow N_1$ is the inclusion map, the map $M_1\to M_0$ is $\mu $, and the map $L_0\twoheadrightarrow \coker \lambda $ is the natural projection. The only maps that need sufficient explaining are the maps $N_1\to M_1 $ and $M_0\to L_0$.

Given that $b=\iota (a)$ for $\iota : \ker \lambda \hookrightarrow N_1$ the inclusion map, we define $c$ as the preimage $c\coloneqq \beta _1^{-1} (b)$. Such a $c$ must exist, as $\beta_1$ must be epic for $1_\bullet $ to be exact by lemma \ref{lma:1}. The question is whether this operation is well-defined; that is, different choices $c_1,c_2\mapsto b$ must map to the same end result if we continue along the sequence. To show this, we must first define the map $L_0\to M_0$.

This map is defined similarly. Since the diagram commutes, $\beta_0(d)=\nu (b)$. However, since $b\in \ker \nu $, we have $\beta_0(d)=0$. It follows that $d\in \ker \beta $, and since $0_\bullet $ is exact, we have that $\im \alpha_0=\ker \beta_0$, and there exists an $e\in L_0$ such that $\alpha_0(e)=d$. Additionally, note that since $0_\bullet $ is exact, $\alpha_0$ is mono by lemma \ref{lma:1}, so this $e$ is unique. Therefore, we can define the map $M_0\to L_0$ such that $d\mapsto \alpha_0^{-1} (d)$, and this is well-defined.

Now to show the map $b\mapsto c$ is well-defined. Let $c_1,c_2\in M_1 $ such that $\beta_1(c_1)=\beta_1(c_2)=b$. Since $\beta_1$ is a homomorphism, $\beta_1(c_1-c_2)=0$. Therefore, $c_1-c_2\in \ker \beta_1$. Since $1_\bullet $ is exact, $c_1-c_2\in \im \alpha_1$. Since $1_\bullet $ is exact, $\alpha_1$ is mono by lemma \ref{lma:1}, so there exists precisely one $g\in L_1 $ such that $\alpha_1(g)=c_1-c_2$. Consider the diagram
\[\begin{tikzcd}[row sep = 2.5em, column sep = 2.5em]
	0\arrow[r, dotted, no head]&g\arrow[d, "\lambda "]\arrow[r, "\alpha_1"]&(c_1-c_2)\arrow[r, "\beta _1"]\arrow[d, dotted, no head, "\mu "]&0\arrow[r, dotted, no head]\arrow[d, dotted, no head]&0\\
	0\arrow[r, dotted, no head]&\lambda (g)\arrow[d]\arrow[r, dotted, no head]&\arrow[r, dotted, no head]\arrow[d, dotted, no head]&\,\\
				   &0\arrow[r, dotted, no head]&\,
\end{tikzcd}.\]
We need only explain $\lambda (g)\to 0$. It suffices to show columns are exact in diagram \eqref{eq:1}, since this would imply $\im \lambda =\ker \pi  $ for $\pi : L_0 \to \coker \lambda $ the canonical projection. It suffices to show much less, actually, but I will just show this broader statement. Consider a complex of the form
\[\begin{tikzcd}[row sep = 2.5em, column sep = 2.5em]
	0\arrow[r]&\ker \lambda \arrow[r, hook, "\iota "]&L_1\arrow[r, "\lambda "]&L_0\arrow[r, two heads, "\pi "]&\coker \lambda \arrow[r]&0 
\end{tikzcd}\]

Since $\iota : \ker \lambda \hookrightarrow L_1$ is injective by the universal property of kernels, we have that $\ker \iota =\left\{ 0 \right\} =\im 0$, and the complex is exact at $\ker \lambda $. Note that $\ker \lambda =\iota (\ker \lambda )=\im \iota $, so the complex is exact at $L_1$. Now note that since $\coker \lambda \cong L_0/\im \lambda $, we have that $\im \lambda \subseteq \ker \pi $. Additionally, $x\in \ker \pi $ implies that $\pi (x)=x+\im \lambda =\im \lambda $. By properties of cosets, it follows that $x\in \im \lambda $. Therefore, the complex is exact at $L_0$. Finally, recall the universal property of cokernels: the cokernel $\coker \lambda $ must be initial with respect to morphisms $\beta : L_0 \to Z $ such that the diagram
\[\begin{tikzcd}[row sep = 2.5em, column sep = 2.5em]
	L_1\arrow[rr, bend left, "0"]\arrow[r, "\lambda "]&L_0\arrow[r, "\beta "]\arrow[d, two heads, "\pi "]&Z\\
				&\coker \lambda \arrow[ru, dotted, "\overline{\beta}"]
\end{tikzcd}\]
commutes. Note that $\pi $ is epic, so $\im \pi = \coker \lambda =\ker 0$, so the complex is exact at $\coker \lambda $. Rows are thus exact complexes, and we have that $\pi (\lambda (g))=0$ since $\im \lambda = \ker \pi $.

Now back to showing $b\mapsto c$ is well-defined. Suppose that $c_2\mapsto d\mapsto e$. Since the diagram commutes, we must have $\mu \alpha_1(g)=\alpha_0\lambda (g)$. We also have that $c_1 = c_2+\left( c_1-c_2 \right) =c_2+\alpha_1(g)$. Since $0_\bullet $ is epic, by lemma \ref{lma:1}, $\alpha_0$ is mono, and thus the only thing mapping to $d$ is $e$, and thus $\mu (c_2)=\alpha_0(e)$ by commutativity. Now note that by commutativity of the diagram, $\alpha_0(e+\lambda (g))=\mu (c_1)=\mu (c_2+c_1-c_2)$:
\[\begin{tikzcd}[row sep = 2.5em, column sep = 2.5em]
	&&c_2+c_1-c_2\arrow[d]\\
	&e+\lambda (g)\arrow[r]&d+\alpha_0(\lambda (g))
\end{tikzcd}\]
Since $\lambda (g)$ dies in $\coker \lambda $, we have that $c_1\mapsto f+0=f$ in diagram \eqref{eq:2}. Therefore, $c_1,c_2$ eventually map to the same thing, and $\delta $ is well-defined overall as the map
\[\begin{tikzcd}[row sep = 2.5em, column sep = 2.5em]
	a\arrow[r, mapsto, "\iota "]&b\arrow[r, mapsto, "\beta _1^{-1} "]&c\arrow[r, mapsto, "\mu "]&d\arrow[r, mapsto, "\alpha_0^{-1} "]&e\arrow[r, mapsto, "\pi "]&f.
\end{tikzcd}\]

Now, we must show the complex
\[\begin{tikzcd}[row sep = 2.5em, column sep = 2.5em]
	0\arrow[r, "0"]&\ker \lambda \arrow[r, "\alpha_1|_{\ker \lambda }"]&\ker \mu \arrow[r, "\beta_1|_{\ker \mu }"]&\ker \nu \arrow[r, "\delta "]&\coker \lambda 
\end{tikzcd}\]
is exact at $\ker \nu $. First, we show $\im \beta_1|_{\ker \mu}\subseteq \ker \delta $. Consider the commutative diagram
\[\begin{tikzcd}[row sep = 2.5em, column sep = 2.5em]
	\ker \mu \arrow[r, "\beta_1"]\arrow[d, "\iota "]&\ker \nu \arrow[d, "\iota "]\\
	M_1\arrow[d, "\mu "]&N_1\arrow[l, "\beta_1^{-1} "]\\
	0
\end{tikzcd}\]
Notice that by the commutativity of this diagram, we have that
\[\begin{tikzcd}[row sep = 2.5em, column sep = 2.5em]
	a\arrow[r, mapsto, "\beta_1"]&\beta_1(a)\arrow[r, mapsto, "\iota "]&\beta_1(a)\arrow[r, mapsto, "\beta_1^{-1}"]&a\arrow[r, mapsto, "\mu "]&0\arrow[r, "\pi \alpha_0^{-1}"]&0
\end{tikzcd}\]
Therefore, $\im \beta_1\subseteq \ker \delta $. Now we show the converse. Let $a\in \ker \delta $. It follows that at some point in the sequence
\[\begin{tikzcd}[row sep = 2.5em, column sep = 2.5em]
	\delta : a\arrow[r, mapsto, "\iota "]&b\arrow[r, mapsto, "\beta_1^{-1} "]&c\arrow[r, mapsto, "\mu "]&d\arrow[r, mapsto, "\alpha_0^{-1} "]&e\arrow[r, mapsto, "\pi "]&f,
\end{tikzcd}\]
at least one of these elements must be 0. Also note that if one element is 0, all subsequent elements must also be 0. We will go through this one at a time.

If $a=0$, then $a\in \im \beta_1|_{\ker \mu }$ because $0\mapsto 0$ and $0\in \ker \mu $. If $b=0$, then $a=0$ because the inclusion $\iota $ is injective. If $c=0$, then $\beta_1(0)=b$ implies that $b=0$. If $d=0$, then $\mu (c)=0$, and thus $c\in \ker \mu $. It follows that $\beta_1(c)\in \im\beta_1|_{\ker \mu }$. Note that
\[\begin{tikzcd}[row sep = 2.5em, column sep = 2.5em]
	\ker \mu \arrow[r]\arrow[d]&\ker \nu \arrow[d]\\
	M_1\arrow[r]&M_0
\end{tikzcd}\]
commutes, and thus $\beta_1(c)=a$. Therefore, $a\in \im \beta_1|_{\ker \mu }$. If $e=0$, then because $\alpha_0$ is injective, $d=0$, and $a\in \im \beta_1|_{\ker \mu }$.

If $f=0$, then note that
\[\begin{tikzcd}[row sep = 2.5em, column sep = 2.5em]
	L_1\arrow[r, "\alpha _1"]\arrow[d, "\lambda "]&M_1\arrow[r, "\beta _1"]\arrow[d, "\mu "]&N_1\arrow[d, "\nu "]\\
	\im \lambda \arrow[r, "\alpha _0"]&\im \mu \arrow[r, "\beta _0"]&\im \nu 
\end{tikzcd}\]
commutes. Note that $f=0$ if and only if $e\in \im \lambda $, and we have that $\beta_0\alpha _0(\lambda )=0$ because rows are exact. Then $\nu (b)=\beta_0\alpha_0(\lambda )$, and thus $b\in \ker \nu $. Great, we already knew that. Note that $e$ has a preimage under $\lambda $, which must map to $c$ by commutativity of the diagram. Since rows are exact, $\beta_1(c)=0$, and thus $b=0$. Therefore, $\ker \delta \subseteq \im \beta_1|_{\ker \mu }$, and the complex is exact at $\ker \nu $.
\end{proof}
\appendix
\section{Definitions and Results Used in the Proof}
\begin{lemma}\label{lma:1}
	Let
	\[\begin{tikzcd}[row sep = 2.5em, column sep = 2.5em]
		0\arrow[r]&M_2\arrow[r, "d_1"]&M_1\arrow[r, "d_0"]&M_0\arrow[r]&0
	\end{tikzcd}\]
	be a short exact chain complex. Then $d_1$ is a monomorphism and $d_0$ is an epimorphism.
\end{lemma}
\begin{definition}[Kernel]\label{dfn:1}
	The kernel $\ker \phi $ of an $R$-module homomorphism $\phi : R \to S$ is the $R$-module which is final with respect to mappings $\alpha : S \to Z$ such that $\phi \alpha =0$. That is, $\ker \phi $ is the $R$-module such that
	\[\begin{tikzcd}[row sep = 2.5em, column sep = 2.5em]
		R\arrow[rr, bend left, "0"]\arrow[dr, dotted]\arrow[r, "\phi "]&S\arrow[r, "\alpha "]&Z\\
									       &\ker \phi \arrow[u, hook]
	\end{tikzcd}\]
	commutes.
\end{definition}

\end{document}
